\subsection{Random Design, Bootstrap, and Extensions}
% Source integration: attached Chapter 14 screenshots, Sections 14.6--14.8 and selected problems.

The fixed-design model conditions on the observed predictor values.  Many observational studies are more naturally viewed as sampling random predictor-response pairs, with the fixed-design analysis recovered after conditioning on the realized design.

\begin{definition} \label{def:random-design-linear-regression-model} \hlt{Random-Design Linear Regression}\\
Let $(Y_i,\xi_i)$, $i=1,\ldots,n$, be independent copies of a pair $(Y,\xi)$, with $\xi\in\R^p$.  A random-design linear regression model specifies
\[
\mathbb{E}(Y\mid \xi=x)=x^t\beta,
\qquad
\operatorname{Var}(Y\mid \xi=x)=\sigma^2.
\]
After observing the design rows $x_i=\xi_i^t$, the conditional model is the fixed-design model $Y=X\beta+e$.
\end{definition}

\begin{proposition} \label{prop:fixed-design-intervals-unconditional-coverage} \hlt{Conditional Coverage Implies Unconditional Coverage}\\
Suppose a fixed-design confidence procedure $C(X,Y)$ satisfies
\[
\mathbb{P}\bigl(\beta_j\in C(X,Y)\mid X\bigr)=1-\alpha
\]
for every realized full-rank design matrix $X$.  Then
\[
\mathbb{P}\bigl(\beta_j\in C(X,Y)\bigr)=1-\alpha.
\]
\begin{proof}
Use the tower property:
\[
\mathbb{P}\bigl(\beta_j\in C(X,Y)\bigr)
=\mathbb{E}\left[\mathbb{P}\bigl(\beta_j\in C(X,Y)\mid X\bigr)\right]
=\mathbb{E}(1-\alpha)=1-\alpha.
\]
\end{proof}
\end{proposition}

\begin{remark} \label{rem:random-design-variance-differs} \hlt{Unconditional Variability Can Differ}\\
The fixed-design estimator remains unbiased under the random-design model by conditioning on $X$ and then averaging.  Its unconditional variance, however, also reflects the randomness of $X$ and generally is not the same matrix obtained by simply freezing $X$.
\end{remark}

\begin{method} \label{meth:regression-pairs-bootstrap-cross-reference} \hlt{Pairs Bootstrap for Random-Design Regression}\\
For random-design regression, resample observed pairs $(Y_i,x_i)$ with replacement, refit the regression in each bootstrap sample, and use the empirical spread of the resulting coefficients or predictions as a standard-error estimate.  This is the regression version of Method \ref{meth:pairs-bootstrap-for-regression-statistics}.
\end{method}

\begin{perspective} \label{per:least-squares-extensions-source-chapter} \hlt{Extensions Beyond Ordinary Least Squares}\\
The source chapter flags several common extensions.
\begin{enumerate}[label=(\roman*)]
\setlength{\itemsep}{0pt}
\item \hlt{Robust regression}: least squares can be highly sensitive to outliers; replacing $\sum_i\widehat{e}_i^2$ by $\sum_i\Psi(\widehat{e}_i)$ gives an $M$-estimation analogue, with $\Psi(u)=u^2$ returning least squares.
\item \hlt{Variable selection}: many applications have more candidate predictors than should be included simultaneously; coefficient estimates and their standard errors depend on the selected model.
\item \hlt{Influence}: high-leverage observations can dominate a fit, especially in multiple regression where influence is harder to see in two-dimensional plots.
\item \hlt{Inverse calibration}: after fitting an instrument response curve, predicting an unknown input from a measured response is an inverse problem and has uncertainty not captured by coefficient standard errors alone.
\end{enumerate}
\end{perspective}

\begin{remark} \label{rem:local-linear-smoothing-regression-bridge} \hlt{Local Linear Smoothing}\\
When a global straight line is too rigid but the relationship is locally linear, local linear smoothing fits small weighted lines around each target point.  Method \ref{meth:local-linear-regression-smoothing} records this nonparametric regression version in its natural owner, extending the broader smoothing setup in Definition \ref{def:nonparametric-regression-smoothing}.
\end{remark}
